\section{Counterarguments}
\label{sec:counter}

The argument has real weaknesses. The strongest objections are set out here
in their own words, with what the evidence in this paper says about each.

\begin{counterpoint}[1\enspace ``Remote places got richer, not poorer'']
True. Absolute living standards have risen across almost all small and remote
states, and nothing here shows otherwise. The claim is relative: about
position, volatility and control. Where the data show relative gaps
narrowing---internet access, shipping relative to the old hubs, diabetes
relative to the world---the paper says so.
\end{counterpoint}

\begin{counterpoint}[2\enspace ``It's governance, not geography'']
Partly. Government effectiveness absorbs about half of the cross-country
income penalty of remoteness in \cref{tab:regression}, and colonial history
and resource endowments are uncontrolled confounders. But distance retains a
large, precisely estimated association with income even after governance is
held fixed---and the direction of causation between geography and governance
plausibly runs both ways, since thin institutions are themselves one of the
taxes of smallness.
\end{counterpoint}

\begin{counterpoint}[3\enspace ``Technology will solve it'']
For access, technology is solving it (\cref{fig:internet}). Satellite internet
and mobile money are genuine equalisers. But redundancy, intensity of use,
price and the capacity to build on each layer are not converging as quickly
(\cref{fig:cables,fig:adoption}), and new technologies create new
dependencies on a small number of foreign providers.
\end{counterpoint}

\begin{counterpoint}[4\enspace ``Small states are not passive victims'']
Correct, and important. Small states earn real rents from their position.
The Parties to the Nauru Agreement, a bloc of Pacific states, transformed their
bargaining power over foreign tuna fleets with the Vessel Day Scheme, which
sells fishing days rather than licences \parencite{aqorau2009}. Several
states run international ship registries. Tuvalu licenses its country-code
domain, .tv, to a global registry operator \parencite{pina2021tv}. These are
evidence of agency---and of the edge being turned into an asset.
\end{counterpoint}

\begin{counterpoint}[5\enspace ``Diet and health have many causes'']
Agreed. The paper does not attribute Pacific diabetes rates to remoteness or
trade. It argues only that remote, import-dependent food systems and thin
health workforces make prevention and treatment more costly.
\end{counterpoint}

\begin{counterpoint}[6\enspace ``Globalisation didn't cause climate change in the Pacific'']
Also agreed. The climate argument is about an asymmetry of exposure versus
influence, not causation by trade.
\end{counterpoint}
